\chapter{Predicción supervisada del movimiento diario}
\label{ch:o4}

El capítulo anterior cerró con un modelo no supervisado que separaba los dos regímenes
de comportamiento de \textit{Larus fuscus}, los días sedentarios de los de migración,
pero que no predecía dónde estaría el ave al día siguiente. Este capítulo retoma esa
pregunta, ya pendiente desde el final del capítulo \ref{ch:o2}, y la responde con
modelos supervisados que combinan tres fuentes de información: la posición diaria del
capítulo \ref{ch:o1}, la cinemática del tramo recién recorrido y la etiqueta de
régimen producida por el HMM del capítulo \ref{ch:o3}. La pregunta concreta es la
misma que en la cadena de Markov visible, dado dónde está hoy el ave, ¿dónde estará
mañana?, pero la respuesta dispone ahora de variables adicionales y de algoritmos no
lineales, y se evalúa contra los dos mismos rivales: la persistencia trivial y la
cadena de Markov de primer orden.

El capítulo plantea dos formulaciones del problema. La primera es la natural por
simetría con el capítulo de Markov, una clasificación de la celda del día siguiente
sobre la misma rejilla de 0,5\textdegree; sirve para caracterizar el techo del
enfoque y diagnosticar dónde falla. La segunda reformula el objetivo como una
regresión continua del desplazamiento diario y añade una banda de incertidumbre
interpretable, que es la salida natural para llevar la predicción a un mapa. Las dos
formulaciones se entrenan sobre las mismas variables y la misma partición temporal y
se comparan contra los mismos baselines, de modo que cualquier diferencia entre ellas
es atribuible al planteamiento del problema y no a la preparación de los datos.

El producto del capítulo, una predicción por ave y día con su banda de incertidumbre,
alimenta los mapas del capítulo \ref{ch:o5}, que cartografían el error y lo
estratifican por régimen y por mes. La preparación común a las dos formulaciones
(datos, variables, familias, métricas y baselines) se fija en la sección
\ref{sec:o4-diseno} y no se vuelve a tocar; las dos formulaciones se desarrollan en
las secciones \ref{sec:o4-categorica} y \ref{sec:o4-cuantiles}, y la
\ref{sec:o4-sintesis} consolida el techo común y deja definido lo que recibe el
capítulo siguiente.

\section{Datos, variables y diseño experimental}
\label{sec:o4-diseno}

\subsection{Entrada y partición temporal}
\label{sec:o4-datos}

Las dos formulaciones del capítulo parten de los mismos dos entregables: la tabla
diaria del capítulo \ref{ch:o1} (una posición por ave y día natural, con los huecos
representados de forma explícita) y la tabla de cinemática y régimen del capítulo
\ref{ch:o3} (la cinemática entrante de cada día, la etiqueta de estado del HMM causal
y su probabilidad asociada, más la columna \texttt{split} que asigna cada fila a
entrenamiento, validación o prueba). La partición temporal del HMM se reutiliza
intacta, sin volver a inducirla, para garantizar que el HMM y los modelos
supervisados se evalúan sobre los mismos días futuros de las mismas aves.

Sobre esa unión se construye el conjunto candidato del capítulo. Cada fila es un par
(ave, día $t$) con su día siguiente $t+1$ disponible: el día $t$ debe tener
cinemática válida del HMM, su posición y la de $t+1$ deben caer en una celda activa
de la rejilla del capítulo \ref{ch:o2}, y el día $t+1$ debe ser el día natural
inmediatamente posterior al $t$ (los pares que atraviesan un hueco calendario se
descartan, para evitar que el modelo aprenda transiciones inexistentes). El resultado
son 20\,188 pares candidatos, repartidos en 14\,598 de entrenamiento, 1\,630 de
validación interna y 3\,960 de prueba, una proporción 72/8/20 que las 82 aves
mantienen en los tres subconjuntos. Sobre estos pares se entrenan y se evalúan todos
los modelos del capítulo, supervisados y baselines, de modo que la comparación es en
igualdad estricta.

La elección de partición es deliberada. La alternativa más exigente, dejar fuera del
entrenamiento todas las observaciones de un ave (\textit{leave-one-bird-out}),
reproduce el fallo diagnosticado en el capítulo \ref{ch:o2}: en ese protocolo, una
fracción importante de los orígenes del conjunto de prueba no estaba representada en
el entrenamiento y la cadena de Markov caía a la distribución marginal del mes. La
partición temporal por ave que se adopta aquí evalúa, en cambio, el escenario
operativo del capítulo \ref{ch:o5}: predecir el futuro próximo de aves cuyo
comportamiento histórico se ha observado en el pasado. En este protocolo, la fracción
de orígenes del conjunto de prueba que el entrenamiento no ha visto baja al 10,7\,\%
(y la de destinos no vistos al 11,1\,\%), lo que sitúa al modelo en un régimen
razonable para juzgar su capacidad predictiva real, sin el sesgo pesimista del
\textit{leave-one-bird-out} ni el optimista de una partición aleatoria.

\subsection{Variables predictivas}
\label{sec:o4-features}

El modelo dispone de diez variables, que se agrupan en tres bloques. La posición
actual del ave aporta el contexto geográfico (\texttt{lat}, \texttt{lon}). El ciclo
anual se codifica con una pareja de variables cíclicas
($\sin\!\tfrac{2\pi\, \mathrm{doy}}{365}$, $\cos\!\tfrac{2\pi\, \mathrm{doy}}{365}$,
donde \textit{doy} es el día del año), que evita la discontinuidad artificial del 31
de diciembre al 1 de enero. El movimiento reciente se resume con la cinemática del
tramo entrante $t\!-\!1 \!\to\! t$ del capítulo \ref{ch:o3}, ya definida allí en sus
términos: el desplazamiento del tramo (\texttt{step\_in\_km}), su dirección codificada
como seno y coseno del rumbo (\texttt{sin\_bearing\_in}, \texttt{cos\_bearing\_in}) y
el coseno del giro respecto al tramo anterior (\texttt{cos\_turning\_in}). Y la salida
del HMM aporta el régimen como dos variables, la etiqueta de estado
(\texttt{state\_b\_causal}) y la probabilidad de migración del día
(\texttt{posterior\_b\_migracion\_causal}), ambas decodificadas por filtrado
\textit{forward-only}.

Todas las variables respetan un mismo criterio de diseño: ninguna observa el día
$t+1$ que se quiere predecir. Es un criterio firme. Una variable que dependa de la
posición de $t+1$ filtra el objetivo y produce métricas optimistas que no
corresponden al uso real del modelo, en el que el día $t+1$ aún no se ha observado.
La cinemática entrante satisface el criterio porque resume el movimiento ya
recorrido hasta el final del día $t$, la inercia real con la que el ave llega al
arranque del día siguiente; la cinemática \emph{saliente} ($t \!\to\! t\!+\!1$), en
cambio, lo violaría. La etiqueta de régimen del HMM también lo satisface por la
construcción causal del capítulo \ref{ch:o3}: el filtrado \textit{forward-only} computa
la probabilidad del estado en $t$ a partir únicamente de las observaciones hasta $t$,
sin mirar al futuro.

\subsection{Familias supervisadas y configuración}
\label{sec:o4-modelos}

El capítulo compara las tres familias de aprendizaje basadas en árboles introducidas
en el capítulo \ref{ch:preliminares}, sección \ref{sec:prelim-ml}: bosques aleatorios
\cite{breiman2001}, XGBoost \cite{chen2016} y LightGBM \cite{ke2017}. Es la
comparativa comprometida en la introducción y la que cierra la simetría con los
modelos no supervisados de los capítulos anteriores: tres algoritmos de capacidad
similar, ajustados al mismo problema, los mismos datos y la misma métrica.

Cada familia se entrena con una configuración fija y conservadora de hiperparámetros,
sin barrido por rejilla. La configuración se resume en la tabla
\ref{tab:o4-hyperparams}. El criterio es metodológico, no de comodidad: un barrido por
rejilla optimiza el rendimiento sobre el conjunto de validación interna, lo que
introduce un sesgo de selección sobre esa misma validación que infla las métricas
reportadas y dificulta separar la mejora real del azar entre configuraciones próximas.
El control del sobreajuste se delega aquí en tres mecanismos independientes: los
hiperparámetros conservadores (profundidad acotada, hojas con soporte mínimo y muestreo
estocástico de filas y columnas en \textit{boosting}), la parada temprana basada en la
validación interna (los modelos de \textit{boosting} fijan un máximo de 1\,000
iteraciones y detienen el entrenamiento si la pérdida en validación no mejora en 50
iteraciones consecutivas) y el diagnóstico explícito de la diferencia entre el error
en entrenamiento y en validación.

\begin{table}[H]
  \centering
  \small
  \begin{tabular}{lll}
    \toprule
    \textbf{Familia} & \textbf{Hiperparámetro} & \textbf{Valor} \\
    \midrule
    \multirow{3}{*}{Bosque aleatorio}
      & \texttt{n\_estimators}    & 300 \\
      & \texttt{max\_depth}       & 12 \\
      & \texttt{min\_samples\_leaf} & 10 \\
    \midrule
    \multirow{6}{*}{XGBoost}
      & \texttt{learning\_rate}   & 0,05 \\
      & \texttt{max\_depth}       & 6 \\
      & \texttt{min\_child\_weight} & 10 \\
      & \texttt{subsample}        & 0,8 \\
      & \texttt{colsample\_bytree} & 0,8 \\
      & \texttt{n\_estimators}    & 1\,000 con parada temprana a 50 \\
    \midrule
    \multirow{5}{*}{LightGBM}
      & \texttt{learning\_rate}   & 0,05 \\
      & \texttt{num\_leaves}      & 31 \\
      & \texttt{min\_child\_samples} & 20 \\
      & \texttt{subsample} / \texttt{colsample\_bytree} & 0,8 / 0,8 \\
      & \texttt{n\_estimators}    & 1\,000 con parada temprana a 50 \\
    \bottomrule
  \end{tabular}
  \caption[Configuración de hiperparámetros]{Configuración fija de hiperparámetros por
  familia, compartida por las dos formulaciones del capítulo. Los valores son
  conservadores y reflejan el consenso de la literatura aplicada para datos tabulares
  de tamaño moderado; el control del sobreajuste se delega en la combinación de la
  configuración, la parada temprana sobre la validación interna y el diagnóstico
  \textit{train-test}.}
  \label{tab:o4-hyperparams}
\end{table}

\subsection{Métricas y baselines}
\label{sec:o4-metricas}

Las métricas de evaluación son las mismas del capítulo \ref{ch:o2}, para que la
comparación entre capítulos sea directa. La proporción de aciertos en el primer
intento (\textit{top-1}) y entre los tres primeros (\textit{top-3}) mide la precisión
puntual sobre la celda; la mediana de la distancia geográfica entre el centroide de la
celda predicha y la posición real cuantifica, en kilómetros, cuánto se desvía cuando
falla. El logaritmo de la pérdida (\textit{log-loss}), heredado del cierre del capítulo de
Markov, evalúa la calidad probabilística completa de la distribución que cada modelo
entrega sobre las celdas, no sólo el acierto del primer intento, y por tanto complementa
a las métricas de precisión puntual con una dimensión que éstas no recogen. La elección
entre familias dentro de cada formulación pondera todas estas métricas en conjunto,
sin que una sola decida por sí sola: una familia con mejor \textit{log-loss} pero
distancia mediana sensiblemente peor o un sobreajuste marcado no se prefiere a una
alternativa más equilibrada. La formulación continua de la sección
\ref{sec:o4-cuantiles} no produce una distribución categórica nativa y se evalúa con
sus propias métricas geométricas allí.

El análisis del error se desagrega en cuatro cortes que se aplican a las dos
formulaciones. El \emph{global} resume el comportamiento sobre todo el conjunto de
prueba. Los cortes por régimen del HMM, \emph{estacionario} y \emph{migración},
distinguen los días en los que el ave permanece en su zona de los de desplazamiento
dirigido (3\,509 frente a 451 días en el conjunto de prueba, según
\texttt{state\_b\_causal}); separar estos dos regímenes es necesario porque una
métrica global está dominada por la moda mayoritaria del objetivo (los días
estacionarios) y oculta el comportamiento del modelo donde el problema es realmente
informativo. El cuarto corte, \emph{días de movimiento real}, selecciona los días en
los que la celda de $t+1$ difiere de la de $t$; es el escenario en el que la
persistencia trivial vale cero por construcción y donde un modelo predictivo tiene
sentido.

Los dos baselines del capítulo \ref{ch:o2} se reutilizan, reentrenados sobre la misma
partición temporal por ave para que la comparación sea en igualdad. La
\emph{persistencia} predice que mañana el ave seguirá en la celda de hoy; su
distribución es degenerada (asigna toda la masa a una celda) y no produce
\textit{log-loss} comparable, pero sigue siendo la referencia exigente que el resto
del trabajo persigue. La \emph{cadena de Markov de primer orden con estacionalidad
mensual}, ajustada sobre el conjunto de entrenamiento de este capítulo, es el rival
con distribución probabilística completa: un modelo no lineal supervisado debe, al
menos, mejorarlo para justificar su mayor complejidad.

\section{Predicción categórica de la celda}
\label{sec:o4-categorica}

La primera formulación es la natural por simetría con el capítulo \ref{ch:o2}: el
objetivo de la predicción es la celda de 0,5\textdegree{} del día siguiente, tratada
como una categoría más entre las 871 celdas activas que las posiciones del conjunto
de entrenamiento ocupan en algún momento. La rejilla y el horizonte son los mismos
que en la cadena de Markov, de modo que el \textit{top-1}, el \textit{top-3} y la
distancia mediana del modelo supervisado son directamente comparables con los de los
dos baselines. La pregunta operativa que esta sección responde no es sólo si el
supervisado bate a los baselines, sino dónde y por qué falla, porque ese diagnóstico
es el que motiva la reformulación de la sección \ref{sec:o4-cuantiles}.

\subsection{Planteamiento del problema}
\label{sec:o4-cat-planteamiento}

Cada par (ave, día~$t$) del conjunto candidato de la sección \ref{sec:o4-datos} se
convierte en un ejemplo de entrenamiento o de evaluación: la entrada son las diez
variables de la sección \ref{sec:o4-features} en el día $t$, la etiqueta es el
identificador de la celda que la posición de $t+1$ ocupa. Las tres familias de la
sección \ref{sec:o4-modelos} se ajustan en modo \emph{poblacional}: un único modelo
sobre las 82 aves, sin la identidad del individuo (\texttt{bird\_id}) entre las
variables. La decisión es deliberada y se sostiene desde dos argumentos. El primero
es de uso: el modelo poblacional cubre tanto a las aves ya observadas en
entrenamiento como a un ave nueva no vista, y no hace falta un modelo distinto por
individuo. El segundo es empírico, y se cierra en la sección
\ref{sec:o4-cat-diagnostico}: un clasificador entrenado únicamente con los datos del
ave de mayor histórico no supera al poblacional evaluado sobre las mismas
observaciones, lo que indica que la distribución poblacional ya captura la
heterogeneidad inter-individual sin necesidad de personalizar.

\subsection{Resultados globales y descarte de LightGBM}
\label{sec:o4-cat-globales}

La comparativa global sobre el conjunto de prueba se resume en la tabla
\ref{tab:o4-cat-globales} y se visualiza en la figura \ref{fig:o4-cat-modelos}. El
bosque aleatorio y XGBoost rinden de forma casi indistinguible entre sí en las
cuatro métricas y quedan en la misma banda; LightGBM, sobre los mismos datos, se
separa visiblemente por uno o dos órdenes de magnitud en \textit{log-loss} y en
distancia mediana, valores incompatibles con los de las otras dos familias.

\begin{table}[H]
  \centering
  \footnotesize
  \begin{tabular}{lrrrr}
    \toprule
    \textbf{Modelo} & \textbf{\textit{top-1}} & \textbf{\textit{top-3}} & \textbf{\textit{log-loss}} & \textbf{Dist.\,mediana (km)}\\
    \midrule
    Persistencia       & 0,778 & 0,778 &       & 20,9 \\
    Cadena de Markov   & 0,552 & 0,637 &       & 24,7 \\
    Bosque aleatorio   & 0,576 & 0,732 & 5,49  & 24,1 \\
    XGBoost            & 0,582 & 0,702 & 5,60  & 24,3 \\
    LightGBM           & 0,359 & 0,579 & 13,12 & 147,7 \\
    \bottomrule
  \end{tabular}
  \caption[Métricas globales sobre el conjunto de prueba]{Métricas globales sobre el
  conjunto de prueba (3\,960 pares (ave, día), el último 20\,\% de los días de cada
  ave). El \textit{log-loss} se computa sobre las 871 celdas activas del
  entrenamiento; los baselines no producen una distribución comparable y se dejan en
  blanco en esa columna. La distancia es la mediana entre el centroide de la celda
  predicha y la posición real. La persistencia domina las cuatro métricas
  comparables y LightGBM queda visiblemente separado por su divergencia en
  validación interna (figura \ref{fig:o4-cat-curvas}).}
  \label{tab:o4-cat-globales}
\end{table}

\begin{figure}[H]
  \centering
  \includegraphics[width=\linewidth]{o4_fig04_models-comparison.png}
  \caption[Comparativa global de los modelos supervisados y los baselines]{Métricas
  globales sobre el conjunto de prueba (3\,960 pares (ave, día), $20\,\%$ último de
  cada ave) de las tres familias supervisadas en modo poblacional (bosque aleatorio,
  XGBoost y LightGBM) frente a los dos baselines del capítulo \ref{ch:o2}
  (persistencia y cadena de Markov de primer orden). Los cinco paneles muestran, de
  izquierda a derecha, el \textit{top-1} y el \textit{top-3} globales, la mediana de
  la distancia en kilómetros, el \textit{log-loss} sobre las 871 celdas activas y el
  \textit{top-1} restringido a los días en los que el ave cambia efectivamente de
  celda. La persistencia domina las métricas globales y LightGBM queda
  visiblemente separado por su \textit{log-loss} y su distancia anómalos, coherentes
  con la divergencia diagnosticada en la figura \ref{fig:o4-cat-curvas}.}
  \label{fig:o4-cat-modelos}
\end{figure}

La separación de LightGBM no es una variabilidad estadística sino un fallo de ajuste
diagnosticable, que la figura \ref{fig:o4-cat-curvas} expone. La curva del
\textit{log-loss} sobre la validación interna empieza ya por encima del valor que
daría una distribución uniforme sobre las 871 celdas ($\log 871 \approx 6{,}77$) y
sube de forma monótona desde la primera iteración hasta estabilizarse en torno a 32,
lo que activa la parada temprana en la iteración 1; XGBoost, sobre los mismos datos y
con un esquema de regularización análogo, desciende como cabe esperar y se detiene en
torno a la iteración 320. La interpretación es que la configuración conservadora de
LightGBM, pensada para evitar el sobreajuste sobre datos tabulares de complejidad
moderada, no extrae señal generalizable de un objetivo de muy alta cardinalidad y
distribución desbalanceada como el del problema categórico aquí planteado: con 871
clases activas en entrenamiento y un peso fuertemente desigual entre ellas, las hojas
de pequeño soporte y la división del crecimiento por hojas que caracterizan a
LightGBM no consiguen aproximar la distribución condicional, mientras que el
crecimiento por niveles y la regularización explícita de XGBoost sí lo hacen. La
familia se descarta como candidato sobre la formulación categórica; en la sección
\ref{sec:o4-cuantiles}, cuando el objetivo deja de ser categórico, vuelve a ser
viable y se reincorpora a la comparativa.

\begin{figure}[H]
  \centering
  \includegraphics[width=\linewidth]{o4_fig03_learning-curves.png}
  \caption[Curvas de aprendizaje sobre la validación interna]{Evolución del
  \textit{log-loss} sobre la validación interna a lo largo de las iteraciones de
  \textit{boosting} para XGBoost (izquierda) y LightGBM (derecha), en modo
  poblacional. La línea vertical punteada marca la iteración en la que la parada
  temprana (paciencia de 50) detiene el entrenamiento. XGBoost desciende
  monótonamente hasta estabilizarse en torno a la iteración 320; LightGBM, en
  cambio, arranca ya por encima del valor de una distribución uniforme sobre las 871
  clases activas ($\log 871 \approx 6{,}77$) y crece sin pausa, lo que indica que su
  configuración conservadora no extrae señal generalizable de este objetivo
  multiclase de alta cardinalidad y distribución desbalanceada.}
  \label{fig:o4-cat-curvas}
\end{figure}

\subsection{Dónde y por qué falla el modelo}
\label{sec:o4-cat-error}

La lectura global de la tabla \ref{tab:o4-cat-globales} es asimétrica según el rival.
Frente a la cadena de Markov, las dos familias supervisadas que quedan ganan en las
tres métricas comparables (\textit{top-1}, \textit{top-3} y distancia mediana), con
una mejora modesta pero consistente que confirma que las variables del HMM y la no
linealidad del modelo aportan algo sobre la transición monoparamétrica de Markov.
Frente a la persistencia, en cambio, los dos modelos supervisados pierden en
\textit{top-1}: la persistencia, de distribución degenerada, asigna toda la masa
a una celda y su \textit{top-3} coincide siempre con su \textit{top-1}; aun así,
supera al mejor supervisado por 20 puntos en \textit{top-1}. En \textit{top-3} la
brecha se estrecha: el bosque aleatorio alcanza 0,732 frente al 0,778 de la
persistencia, una diferencia de 5 puntos frente a los 20 de \textit{top-1}.
El \textit{log-loss} del mejor supervisado (5,49) es
informativo en términos absolutos (mejor que el 6,77 de una distribución uniforme
sobre las 871 clases activas), pero la persistencia no produce distribución
comparable y no participa en esa columna.

La asimetría tiene una causa estructural directa: en el conjunto de prueba el ave
permanece en la misma celda el 77,8\,\% de los días, una moda absoluta del objetivo
que la persistencia explota por construcción y que cualquier modelo que asigne masa a
un movimiento paga proporcionalmente en \textit{top-1}. Es la geometría del problema
antes que el algoritmo. El desglose por régimen del HMM y por movimiento real, en la
tabla \ref{tab:o4-cat-estado} y en la figura \ref{fig:o4-cat-estado}, refina esta
lectura.

\begin{table}[H]
  \centering
  \footnotesize
  \begin{tabular}{lcccccc}
    \toprule
     & \multicolumn{2}{c}{\textbf{Estacionario}} & \multicolumn{2}{c}{\textbf{Migración}} & \multicolumn{2}{c}{\textbf{Movim.~real}}\\
     & \multicolumn{2}{c}{$n=3\,509$} & \multicolumn{2}{c}{$n=451$} & \multicolumn{2}{c}{$n=880$}\\
    \cmidrule(lr){2-3}\cmidrule(lr){4-5}\cmidrule(lr){6-7}
    \textbf{Modelo} & \textit{top-1} & \textit{top-3} & \textit{top-1} & \textit{top-3} & \textit{top-1} & \textit{top-3}\\
    \midrule
    Persistencia       & 0,841 & 0,841 & 0,284 & 0,284 & 0,000 & 0,000 \\
    Cadena de Markov   & 0,607 & 0,694 & 0,120 & 0,195 & 0,089 & 0,373 \\
    Bosque aleatorio   & 0,632 & 0,792 & 0,137 & 0,259 & 0,164 & 0,445 \\
    XGBoost            & 0,639 & 0,762 & 0,140 & 0,235 & 0,175 & 0,445 \\
    \bottomrule
  \end{tabular}
  \caption[Acierto \textit{top-1} y \textit{top-3} desagregados por régimen del HMM y por movimiento
  real]{Acierto (\textit{top-1} y \textit{top-3}) sobre el
  conjunto de prueba, desagregados por el régimen del HMM causal (\emph{estacionario}
  y \emph{migración}, según \texttt{state\_b\_causal} del capítulo \ref{ch:o3}) y
  restringidos a los días en que la celda de $t+1$ difiere de la de $t$
  (\emph{movim.~real}). Se omite LightGBM, descartado en la sección anterior.
  La persistencia, de distribución degenerada, iguala siempre su \textit{top-1} y
  su \textit{top-3}. Los modelos supervisados reducen notablemente la brecha con
  la persistencia al pasar a \textit{top-3}, y en los días de movimiento real
  la superan con claridad (0,445 frente a 0,000).}
  \label{tab:o4-cat-estado}
\end{table}

\begin{figure}[H]
  \centering
  \includegraphics[width=\linewidth]{o4_fig07_error-by-state-poblacional.png}
  \caption[\textit{Top-1} y \textit{top-3} desagregados por régimen del HMM]{Acierto
  (\textit{top-1} y \textit{top-3}) del modelo poblacional (XGBoost, derecha) y del
  modelo individual del ave 91916A (bosque aleatorio, izquierda) desagregados por el
  régimen del HMM causal del capítulo \ref{ch:o3}, en el conjunto de prueba. El
  régimen \emph{estacionario} cubre 3\,509 días (88,6\,\% del conjunto) y el régimen
  \emph{migración}, 451 días (11,4\,\%). Para el poblacional, el \textit{top-1} cae de
  0,64 en estacionario a 0,14 en migración; el \textit{top-3} cae de 0,76 a 0,24.
  La brecha con la persistencia es mayor en \textit{top-1} que en \textit{top-3}, y
  se concentra en los días de migración (tabla \ref{tab:o4-cat-estado}).}
  \label{fig:o4-cat-estado}
\end{figure}

Dos lecturas. La primera, la propia caída en migración: el \textit{top-1} de los
supervisados se desploma de 0,64 a 0,14 entre los dos regímenes, y la distancia
mediana cuando fallan se dispara de 23 a más de 150\,km. Los días informativos para
la predicción son la minoría y son precisamente donde el modelo pierde casi toda su
capacidad predictiva. La segunda, que la persistencia bate al supervisado también
dentro del régimen de migración del HMM, podría parecer contradictoria con la
motivación de incorporar el régimen como variable, pero es coherente con la
definición causal del estado: la etiqueta \texttt{state\_b\_causal} se asigna a
partir del paso \emph{entrante} ($t\!-\!1 \to t$), de modo que un día puede
pertenecer al régimen de migración por haber llegado tras un desplazamiento grande,
sin que el ave haya iniciado otro desplazamiento durante $t$ ni vaya a iniciarlo en
$t+1$. En esos días, una fracción importante de los 451 marcados como migración, la
persistencia sigue acertando.

El espacio donde un modelo predictivo aporta algo que la persistencia no puede dar
por construcción es, por tanto, el de los \emph{días de movimiento real}: los 880
días del conjunto de prueba en los que la celda de $t+1$ difiere de la de $t$
(el 22,2\,\% del total). Allí la persistencia acierta cero en \textit{top-1} y
en \textit{top-3} por construcción; los modelos supervisados aciertan entre 0,16 y
0,18 en \textit{top-1} y llegan a 0,445 en \textit{top-3}, frente al 0,373 de la
cadena de Markov. El margen es modesto en \textit{top-1}, pero en \textit{top-3}
los supervisados superan claramente a Markov y aplican sobre la persistencia una
ventaja que este contexto permite cuantificar; es el único corte en el que ambas
comparaciones resultan favorables al modelo supervisado.

\subsection{De dónde sale la señal}
\label{sec:o4-cat-importancia}

La descomposición por importancia relativa de las variables, en la figura
\ref{fig:o4-cat-importancia}, ordena cuantitativamente las tres fuentes de
información de las que dispone el modelo. La posición acapara la mayor parte: la
latitud y la longitud suman 0,75 de la importancia normalizada en el modelo
poblacional, lo que reproduce numéricamente la dominación del régimen estacionario
detectada en la sección anterior. El bloque de variables del HMM contribuye 0,05
(la probabilidad de migración) más 0,03 (la etiqueta de estado), una suma modesta
pero coherente: el aporte del HMM al modelo supervisado pasa principalmente por la
\emph{probabilidad blanda} del régimen, que matiza el comportamiento del clasificador
en la franja ambigua entre los días claramente estacionarios y los claramente
migratorios, mientras que la \emph{etiqueta dura} del estado, que en su mayor parte
sólo replica información ya recogida en la probabilidad, pesa poco más de la mitad que
aquella y no llega a ser determinante. El bloque del ciclo anual (seno y coseno del
día del año) suma 0,09 y la cinemática entrante (desplazamiento, rumbo y giro)
alrededor de 0,09. Las variables del HMM, por tanto, están en el mismo orden de
magnitud que el ciclo anual y la cinemática: aportan señal real al supervisado,
modesta y principalmente a través del posterior blando.

\begin{figure}[H]
  \centering
  \includegraphics[width=\linewidth]{o4_fig09_feature-importance-winners.png}
  \caption[Importancia relativa de las variables]{Importancia relativa de las diez
  variables predictivas para el modelo poblacional (XGBoost, izquierda) y para el
  bosque aleatorio entrenado únicamente con el ave 91916A (derecha; ver
  \ref{sec:o4-cat-diagnostico}). En ambos casos la posición (latitud y longitud)
  acapara la mayor parte de la importancia (0,75 en el poblacional, 0,63 en el
  modelo individual). Las variables del HMM aportan 0,05 (probabilidad de migración)
  más 0,03 (etiqueta de estado) en el modelo poblacional y suben a 0,13 más 0,004
  en el modelo individual, donde la señal geográfica se reduce y el régimen blando
  pasa a tener un peso mayor. La etiqueta dura del estado pesa siempre menos que la
  probabilidad blanda; el aporte del HMM al supervisado pasa principalmente por esta última.}
  \label{fig:o4-cat-importancia}
\end{figure}

\subsection{Diagnóstico del techo y cierre del alcance}
\label{sec:o4-cat-diagnostico}

La formulación categórica permite, por tanto, situar dos causas estructurales del
techo de rendimiento que las dos secciones siguientes recogen. La primera es que el
objetivo categórico desperdicia la geometría del espacio: un error que predice una
celda vecina (un desplazamiento de unas decenas de kilómetros) se penaliza
exactamente igual que un error que predice una celda en el otro extremo del dominio
(varios miles), y el modelo no recibe ninguna información sobre la distancia entre
clases en su función de pérdida. Bajo esta restricción, los días de movimiento real
quedan condenados a un \textit{top-1} bajo, porque acertar la celda exacta de
destino es un problema más estricto que predecir la dirección y la magnitud del
desplazamiento. La sección \ref{sec:o4-cuantiles} reformula el objetivo justamente
para devolverle al modelo la información geométrica que la formulación categórica
descarta. La segunda causa es de horizonte y de variables: a un día con variables
locales (posición, cinemática del tramo entrante, ciclo anual y régimen del HMM) le
falta el predictor causal de un evento migratorio (la disponibilidad de vientos
favorables, la historia multi-día del desplazamiento, el destino al que el ave se
dirige), y ninguna combinación de algoritmo y arquitectura puede inventar lo que las
variables no contienen. Esta segunda causa es común a las dos formulaciones del
capítulo y se argumenta de forma cuantitativa, ya con los resultados de ambas a la
vista, en la sección \ref{sec:o4-sintesis}.

Un último apunte sobre el alcance del entrenamiento cierra esta sección y deja
zanjada la elección poblacional anunciada en la sección
\ref{sec:o4-cat-planteamiento}. Un bosque aleatorio entrenado únicamente con los
1\,458 días de entrenamiento del ave de mayor histórico (91916A) y evaluado sobre
sus 406 días de prueba alcanza un \textit{top-1} de 0,660 y un \textit{log-loss} de
4,33; el modelo poblacional restringido a las mismas 406 observaciones obtiene
\textit{top-1} 0,653 y \textit{log-loss} 2,72, prácticamente igual en acierto y
sensiblemente mejor en calidad probabilística. La heterogeneidad inter-individual
que el capítulo \ref{ch:o1} anticipaba y el capítulo \ref{ch:o3} cuantificaba existe,
pero la distribución poblacional ya la captura tan bien como un modelo dedicado al
individuo con más datos posible, sin la complejidad ni la falta de escalabilidad de
entrenar un modelo por ave. La formulación continua de la sección siguiente confirma
este resultado con un experimento análogo y diferente, y la sección
\ref{sec:o4-sintesis} consolida el cierre.

\section{Predicción continua del desplazamiento con cuantiles}
\label{sec:o4-cuantiles}

Esta sección desarrolla una respuesta de modelado a la primera de las dos causas
del techo diagnosticadas en la sección \ref{sec:o4-cat-diagnostico}, la que es
salvable sin cambiar el horizonte ni las variables: reformular el problema como una
regresión del desplazamiento diario en lugar de una clasificación de la celda. Con
las mismas variables, la misma partición temporal y los mismos hiperparámetros
conservadores, la reformulación devuelve al modelo la información geométrica del
espacio que la formulación categórica descartaba (el acierto \textit{top-1} sube 18
puntos a nivel global y la distancia mediana del error en los días de migración del
HMM se reduce a un tercio) y, al estimarse en forma de cuantiles, entrega además
una predicción que sale del modelo como un punto del plano acompañado de un
intervalo $[p_{10}, p_{90}]$ por componente, una banda de incertidumbre
interpretable que la clasificación no podía dar y que el capítulo \ref{ch:o5} lleva
al mapa.

\subsection{Reformulación del objetivo y arquitectura del modelo}
\label{sec:o4-q-planteamiento}

El objetivo de la predicción deja de ser el identificador discreto de la celda y
pasa a ser el desplazamiento del ave entre dos días consecutivos, expresado en
grados: $\Delta\mathrm{lat} = \mathrm{lat}(t\!+\!1) - \mathrm{lat}(t)$ y
$\Delta\mathrm{lon} = \mathrm{lon}(t\!+\!1) - \mathrm{lon}(t)$. La elección de la
parametrización cartesiana frente a la polar (distancia, rumbo) es deliberada y
evita la discontinuidad angular en $\pm 180^\circ$, que rompería la suavidad del
objetivo para los pasos de orientación este-oeste y obligaría a un tratamiento
especial. Las dos componentes se predicen por separado y de forma independiente, con
un modelo por componente.

Para acompañar cada predicción puntual de una banda de incertidumbre se ajustan, para
cada componente, tres cuantiles condicionales en lugar de la media: $p_{10}$, $p_{50}$
y $p_{90}$. La estimación minimiza la pérdida cuantílica
(capítulo \ref{ch:preliminares}, sección \ref{sec:prelim-cuantil},
ecuación \ref{eq:prelim-pinball}), que penaliza de forma asimétrica los errores por
exceso y por defecto según el cuantil buscado y, en su agregado sobre los tres
cuantiles, permite ajustar simultáneamente la mediana condicional $p_{50}$ (la
predicción puntual del desplazamiento) y los dos extremos del intervalo
$[p_{10}, p_{90}]$ (una banda de incertidumbre con cobertura nominal del 80\,\%). La
monotonía $p_{10} \leq p_{50} \leq p_{90}$ no la garantiza el ajuste cuando cada
cuantil se estima con un modelo independiente y se fuerza, cuando es necesario, por
ordenación a posteriori; la fracción de filas afectadas y las diferencias entre
familias se analizan en la sección \ref{sec:o4-q-incertidumbre}.

Las tres familias del capítulo se comparan también en esta formulación. XGBoost y
LightGBM minimizan directamente la pérdida cuantílica para cada cuantil, mientras
que el bosque aleatorio, cuyos árboles parten por reducción de la varianza y no
optimizan directamente esa pérdida, obtiene los cuantiles \emph{a posteriori} como
cuantiles empíricos de la distribución de los valores del objetivo en cada hoja, en
la formulación de bosque de regresión cuantílica de \cite{meinshausen2006}. La
ablación per-individuo, que en la formulación categórica se ensayó con un bosque
aleatorio entrenado únicamente con el ave de mayor histórico, se repite aquí con
XGBoost y se comenta en la sección \ref{sec:o4-q-cierre}.

Para comparar con la formulación categórica de la sección \ref{sec:o4-categorica} y
con los baselines, la predicción puntual $p_{50}$ se discretiza al \textit{grid} de
0,5\textdegree{} (un \textit{top-1 mapeado} que cuenta como acierto las
predicciones cuyo punto cae en la celda real) y la distancia se reporta tanto
\emph{nativa} (la distancia entre el punto predicho y la posición real, sin
discretizar) como \emph{vía centroide} (entre el centroide de la celda en la que cae
el punto predicho y la posición real, comparable con la distancia mediana de la
sección \ref{sec:o4-categorica}).

\subsection{La mejora respecto a la formulación categórica}
\label{sec:o4-q-precision}

La tabla \ref{tab:o4-q-precision} pone lado a lado el mejor modelo de la sección
\ref{sec:o4-categorica} (XGBoost categórico), las tres familias de la regresión de
cuantiles y la persistencia trivial sobre el conjunto de prueba, desagregados por
régimen del HMM y por movimiento real. La figura \ref{fig:o4-q-distancia} visualiza
la distribución completa de la distancia entre el punto predicho y la posición
real para las tres familias de cuantiles.

\begin{table}[H]
  \centering
  \footnotesize
  \begin{tabular}{lcccccccc}
    \toprule
     & \multicolumn{2}{c}{\textbf{Global}} & \multicolumn{2}{c}{\textbf{Estacionario}} & \multicolumn{2}{c}{\textbf{Migración}} & \multicolumn{2}{c}{\textbf{Movim.~real}}\\
     & \multicolumn{2}{c}{$n=3\,960$} & \multicolumn{2}{c}{$n=3\,509$} & \multicolumn{2}{c}{$n=451$} & \multicolumn{2}{c}{$n=880$}\\
    \cmidrule(lr){2-3}\cmidrule(lr){4-5}\cmidrule(lr){6-7}\cmidrule(lr){8-9}
    \textbf{Modelo} & \textit{top-1} & \textit{top-3} & \textit{top-1} & \textit{top-3} & \textit{top-1} & \textit{top-3} & \textit{top-1} & \textit{top-3}\\
    \midrule
    \multicolumn{9}{l}{\emph{Formulación categórica (sección \ref{sec:o4-categorica}, referencia)}}\\
    XGBoost categórico  & 0,582 & 0,702 & 0,639 & 0,762 & 0,140 & 0,235 & 0,175 & 0,445 \\
    \midrule
    \multicolumn{9}{l}{\emph{Formulación continua (esta sección)}}\\
    Bosque aleatorio    & 0,771 & 0,898 & 0,842 & 0,946 & 0,217 & 0,523 & 0,049 & 0,566 \\
    XGBoost             & 0,763 & 0,901 & 0,833 & 0,946 & 0,217 & 0,552 & 0,100 & 0,582 \\
    LightGBM            & 0,762 & 0,900 & 0,835 & 0,946 & 0,197 & 0,543 & 0,085 & 0,578 \\
    \midrule
    \multicolumn{9}{l}{\emph{Baseline}}\\
    Persistencia        & 0,778 & 0,778 & 0,841 & 0,841 & 0,284 & 0,284 & 0,000 & 0,000 \\
    \bottomrule
  \end{tabular}
  \caption[Comparativa entre las dos formulaciones del capítulo y los baselines]{Acierto
  (\textit{top-1} y \textit{top-3}, mapeados al \textit{grid} 0,5\textdegree) sobre el
  conjunto de prueba, para el mejor modelo de la formulación categórica, las tres familias
  de la regresión de cuantiles (modo poblacional) y la persistencia trivial. La predicción
  puntual de la formulación continua es el $p_{50}$ de cada componente, llevado al
  \textit{grid} 0,5\textdegree{} para comparar con la formulación categórica. La
  reformulación sube el \textit{top-1} global de 0,58 a 0,76--0,77 y el \textit{top-3}
  de 0,70 a 0,90, superando en esta última métrica a la persistencia (0,778). En los días
  de migración del HMM, el \textit{top-3} pasa de 0,24 a 0,55, el doble del baseline
  trivial. Entre familias, las tres son prácticamente indistinguibles en los cuatro cortes.}
  \label{tab:o4-q-precision}
\end{table}

\begin{figure}[H]
  \centering
  \includegraphics[width=\linewidth]{o4_fig26_distance-distribution.png}
  \caption[Distribución de la distancia entre el punto predicho y la posición
  real]{Distribución de la distancia vía centroide entre la predicción puntual
  ($p_{50}$) y la posición real, sobre el conjunto de prueba, para las tres
  familias de la regresión de cuantiles y la persistencia trivial (modo
  poblacional). Las distribuciones coinciden en el grueso (la masa concentrada por
  debajo de los 50\,km, atribuible a los días estacionarios) y se separan en la cola
  (los días de movimiento real). La diferencia entre las tres familias supervisadas
  es pequeña a lo largo de toda la distribución.}
  \label{fig:o4-q-distancia}
\end{figure}

El primer mensaje de la tabla es la magnitud de la mejora respecto a la formulación
categórica. La sola reformulación del objetivo, con las mismas variables, el mismo
\textit{split} y los mismos hiperparámetros, levanta el \textit{top-1} global de
0,58 a 0,76 (18 puntos absolutos, un 31\,\% relativo). En \textit{top-3} la mejora
es más pronunciada aún: de 0,70 a 0,90 globalmente, con lo que la formulación
continua supera ya a la persistencia (0,778) en esta métrica. En los días de
migración del HMM, el \textit{top-3} pasa de 0,24 a 0,55, el doble del nivel de
la persistencia (0,28). La mejora se reproduce en los días estacionarios
(\textit{top-1} de 0,64 a 0,83--0,84; \textit{top-3} de 0,76 a 0,946). El
mecanismo es geométrico: como el desplazamiento real es
\emph{nulo} el 77,8\,\% de los días de prueba, el regresor aprende a predecir un
desplazamiento próximo a cero en esos días, su $p_{50}$ recae en la celda de
origen y reproduce el acierto de los \textit{self-loops} que la formulación
categórica no fijaba con su argmax sobre 871 clases. Más en general, la pérdida
cuantílica penaliza el error en distancia (no por igual a todas las clases
incorrectas) y empuja al modelo hacia predicciones geométricamente cercanas, que es
exactamente lo que el problema requiere. La única métrica que se mueve a la baja en
la transición es el \textit{top-1} del corte de movimiento real (de 0,18 a 0,05--
0,10), una contrapartida coherente con el mismo mecanismo (el regresor «se atreve»
menos a apostar fuera de la celda de origen), que la ganancia en \textit{top-3}
(0,58 frente a 0,445 de la formulación categórica) compensa con holgura.

El segundo mensaje es la convergencia entre familias. Sobre el conjunto de prueba el
\textit{top-1} global queda entre 0,762 (LightGBM) y 0,771 (bosque aleatorio), una
banda de menos de un punto, con la distancia mediana indistinguible (20,9\,km en
las tres). El mismo patrón se repite en los cuatro cortes de la
tabla \ref{tab:o4-q-precision} y en la figura \ref{fig:o4-q-distancia}, donde las
tres distribuciones se superponen tanto en el grueso como en la cola. Bajo la
formulación continua, la elección de familia es secundaria para la precisión: tres
algoritmos con mecanismos distintos (\textit{bagging} no paramétrico en el bosque
aleatorio, \textit{boosting} de árboles por niveles en XGBoost y \textit{boosting}
por hojas en LightGBM) llegan al mismo lugar.

Como consecuencia secundaria de la convergencia, LightGBM no diverge aquí. La misma
familia que en la sección \ref{sec:o4-cat-globales} se separó del resto y se
descartó como candidato sobre el problema multiclase de 871 clases activas, sobre
la regresión de cuantiles (con seis modelos, uno por cuantil y eje, en lugar de un
único multiclase de muy alta cardinalidad) se ajusta sin patología y rinde
indistinguiblemente del bosque aleatorio y de XGBoost. La división del crecimiento
por hojas y el soporte mínimo conservador no extraen señal generalizable de un
objetivo de 871 clases con peso fuertemente desigual, pero sí de seis regresiones
cuantílicas escalares cada una. La familia se reincorpora a la comparativa de esta
sección en igualdad con las otras dos, y la elección del modelo recomendado para
los mapas se argumenta en la sección \ref{sec:o4-q-cierre} desde la calidad
probabilística, no desde la precisión.

Frente a la persistencia trivial, la formulación continua se sitúa en una posición
muy distinta a la categórica. Aquélla perdía en todas las métricas globales por un
margen amplio; ésta queda apenas una o dos décimas por debajo en \textit{top-1}
global y, en \textit{top-3}, supera ya a la persistencia (0,90 frente a 0,778).
En los días de migración del HMM, el \textit{top-3} de la formulación continua
(0,55) dobla al de la persistencia (0,28): la formulación aprende que en migración
el ave puede desplazarse a cualquiera de varias celdas vecinas, y distribuye masa
sobre ese vecindario de forma que las tres candidatas más probables ya superan
claramente al único destino que la persistencia propone. La persistencia conserva
una pequeña ventaja en \textit{top-1} por la moda dominante de los \textit{self-loops},
pero la diferencia deja de ser cualitativa: la formulación continua iguala o supera
a la referencia exigente del capítulo según la métrica y el corte, mientras añade
una dimensión que la persistencia no entrega, la banda de incertidumbre que la
sección \ref{sec:o4-q-incertidumbre} caracteriza.

\subsection{Banda de incertidumbre y monotonía de los cuantiles}
\label{sec:o4-q-incertidumbre}

La precisión de la predicción puntual es sólo una de las dimensiones que la
formulación de cuantiles entrega. La otra, ausente por completo en la formulación
categórica, es la \emph{incertidumbre calibrada}: para cada predicción, los
extremos $p_{10}$ y $p_{90}$ delimitan un intervalo que se espera contenga la
posición real el 80\,\% del tiempo si la calibración es correcta. La tabla
\ref{tab:o4-q-incertidumbre} resume las tres métricas que evalúan esta dimensión y
la figura \ref{fig:o4-q-calibracion} la visualiza para la cobertura.

\begin{table}[H]
  \centering
  \footnotesize
  \begin{tabular}{lcccc}
    \toprule
    \textbf{Modelo} & \textbf{Cobertura lat / lon} & \textbf{\textit{Pinball} lat / lon} & \textbf{Cruces de cuantil} \\
    \midrule
    Bosque aleatorio  & 0,862 / 0,859 & 0,081 / 0,062 & 0 \\
    XGBoost           & 0,849 / 0,840 & 0,081 / 0,063 & 54 \\
    LightGBM          & 0,836 / 0,788 & 0,079 / 0,062 & 202 \\
    \bottomrule
  \end{tabular}
  \caption[Calidad probabilística de las tres familias]{Métricas de calidad
  probabilística sobre el conjunto de prueba (3\,960 pares, modo poblacional).
  \emph{Cobertura}: fracción empírica de filas en las que la posición real cae
  dentro del intervalo $[p_{10}, p_{90}]$ predicho, por eje (objetivo nominal del
  80\,\%). \emph{Pinball}: pérdida cuantílica promedio sobre los tres cuantiles, por
  eje. \emph{Cruces de cuantil}: número de filas en las que la ordenación
  $p_{10} \leq p_{50} \leq p_{90}$ no se respetaba antes de la corrección a
  posteriori. Las tres familias tienen una \textit{pinball} prácticamente igual; se
  diferencian en la calibración (LightGBM es el más cercano al nominal) y en la
  coherencia interna de los cuantiles (el bosque aleatorio no cruza por
  construcción).}
  \label{tab:o4-q-incertidumbre}
\end{table}

\begin{figure}[H]
  \centering
  \includegraphics[width=\linewidth]{o4_fig25_calibration-coverage.png}
  \caption[Cobertura empírica del intervalo de predicción]{Cobertura empírica del
  intervalo $[p_{10}, p_{90}]$ frente al objetivo nominal del 80\,\%, por familia y
  por eje, sobre el conjunto de prueba en modo poblacional. Una cobertura próxima al
  80\,\% indica una banda de incertidumbre bien calibrada. LightGBM queda más cerca
  del nominal (con una ligera infracobertura en longitud), XGBoost se sitúa unos
  puntos por encima y el bosque aleatorio sobrecubre, con bandas más anchas en
  ambas direcciones.}
  \label{fig:o4-q-calibracion}
\end{figure}

Las tres familias comparten una pérdida cuantílica prácticamente idéntica (en torno
a 0,08 para la latitud y 0,06 para la longitud), lo que confirma que el ajuste de
los tres cuantiles es de calidad equivalente en la métrica que el problema
optimiza, pero se separan en las dos métricas que la métrica de ajuste no captura
directamente: la calibración empírica del intervalo y la coherencia interna de los
cuantiles. LightGBM queda más cerca del 80\,\% nominal (0,84 y 0,79), XGBoost se
sitúa unos puntos por encima (0,85 y 0,84) y el bosque aleatorio sobrecubre (0,86 y
0,86), con bandas más anchas que las otras dos. La interpretación es directa: el
bosque aleatorio entrega una incertidumbre conservadora, en el sentido de que su
intervalo contiene la posición real con más frecuencia de la que dice, mientras que
LightGBM la entrega más ajustada al nivel anunciado.

La columna de cruces de cuantil refleja una diferencia estructural entre las
familias, no un detalle de implementación. XGBoost y LightGBM ajustan un modelo
independiente por cuantil y por eje, y nada en su entrenamiento garantiza que la
ordenación $p_{10} \leq p_{50} \leq p_{90}$ se respete para todas las observaciones,
de modo que aparecen 54 cruces en XGBoost y 202 en LightGBM, corregibles a
posteriori por simple ordenación pero no eliminables \emph{a priori}. El bosque
aleatorio de cuantiles, en cambio, extrae los tres cuantiles de la misma
distribución empírica de cada hoja, y la ordenación se respeta por construcción:
cero cruces. Es una propiedad estructural del \textit{bagging} no paramétrico en
este problema, no una mejora de ajuste.

\subsection{Balance, alcance del entrenamiento y modelo recomendado}
\label{sec:o4-q-cierre}

El balance de la formulación continua es claro respecto a las dos referencias del
capítulo. Frente a la formulación categórica de la sección \ref{sec:o4-categorica},
la mejora es sustancial y se reproduce en todos los cortes: el \textit{top-1} global
sube 18 puntos y el \textit{top-3} global sube 20 puntos (de 0,70 a 0,90). Frente
a la persistencia trivial, la formulación continua revierte por completo la pérdida
amplia que la formulación categórica acumulaba: en \textit{top-1} queda una o dos
décimas por debajo, y en \textit{top-3} la supera (0,90 frente a 0,778 global;
0,55 frente a 0,28 en los días de migración del HMM). A ello añade la dimensión
que la persistencia no puede dar: una banda de incertidumbre calibrada $[p_{10},
p_{90}]$ con cobertura empírica próxima al 80\,\% nominal
(sección \ref{sec:o4-q-incertidumbre}). Sobre los mapas del capítulo
\ref{ch:o5}, esa banda se traduce en una región geográfica interpretable que
ninguna de las alternativas anteriores podía producir.

La elección del modelo recomendado para los mapas del capítulo siguiente pondera
las tres dimensiones que la tabla \ref{tab:o4-q-precision} y la tabla
\ref{tab:o4-q-incertidumbre} resumen, sin que ninguna decida en solitario. En
precisión las tres familias son indistinguibles, así que la decisión se traslada a
la calidad probabilística: LightGBM es el más cercano al 80\,\% nominal de
cobertura, XGBoost queda como alternativa equivalente con un margen pequeño de
sobrecobertura, y el bosque aleatorio entrega bandas perceptiblemente más anchas
con una garantía estructural de monotonía de los cuantiles. Para los mapas del
capítulo \ref{ch:o5}, el modelo recomendado es la regresión de cuantiles con
LightGBM, por entregar la banda mejor calibrada al nivel anunciado, con XGBoost
como alternativa equivalente.

Un último apunte sobre el alcance del entrenamiento, paralelo al que cierra la
sección \ref{sec:o4-cat-diagnostico}, confirma la decisión poblacional. Un regresor
de cuantiles de XGBoost entrenado únicamente con los 1\,458 días de entrenamiento
del ave de mayor histórico (91916A) y evaluado sobre sus 406 días de prueba alcanza
un \textit{top-1 mapeado} de 0,685 y una distancia mediana vía centroide de
14,6\,km; el modelo poblacional restringido a las mismas 406 observaciones obtiene
\textit{top-1} 0,687 y distancia mediana 14,4\,km, indistinguible del individual en
ambas métricas. Dos experimentos con diseños distintos (una clasificación categórica
en la sección \ref{sec:o4-cat-diagnostico} y una regresión de cuantiles aquí)
convergen, por tanto, en la misma conclusión: la distribución poblacional ya captura
al individuo más favorable tan bien como un modelo dedicado, y la heterogeneidad
inter-individual del conjunto, real, no justifica el coste de entrenar un modelo por
ave ni la pérdida de escalabilidad a las 80 aves con menos histórico. El modelo
canónico del capítulo, en sus dos formulaciones, es el poblacional.

\section{Síntesis y límites del enfoque}
\label{sec:o4-sintesis}

El capítulo ha planteado dos formulaciones del mismo problema de predicción del
movimiento diario y las ha comparado contra los dos baselines heredados del
capítulo \ref{ch:o2}. La formulación categórica caracteriza el techo del enfoque y
diagnostica sus causas; la formulación continua corrige la limitación geométrica de
la primera y entrega, además, una banda de incertidumbre interpretable. En las dos,
el modelo canónico es el poblacional: las pruebas con un clasificador y un regresor
entrenados únicamente con el ave de mayor histórico convergen en la misma
conclusión y cierran la decisión global vs per-individuo que el capítulo
\ref{ch:o1} dejó sembrada y el capítulo \ref{ch:o3} arrastró.

Las dos formulaciones convergen también en un mismo límite, que el capítulo no
puede salvar y que merece nombrarse con precisión: la predicción a un horizonte de
un día con variables locales no contiene el predictor causal de un evento
migratorio (la disponibilidad de vientos favorables, la historia multi-día del
desplazamiento, el destino al que el ave se dirige). No es una limitación del
algoritmo (la convergencia entre las tres familias sobre la regresión de cuantiles
de la sección \ref{sec:o4-q-precision} lo confirma) ni de la formulación del
objetivo (la regresión continua recupera ya la información geométrica que la
clasificación descartaba): es del planteamiento del problema a esta escala. El
trabajo futuro que se deriva es el de incorporar esas tres piezas (las variables de
viento ECMWF descartadas en el capítulo \ref{ch:o1}, la historia multi-día como
secuencia de entrada y una noción explícita del destino), todas fuera del alcance
temporal del trabajo.

El producto que el capítulo entrega al capítulo \ref{ch:o5} es la predicción del
modelo recomendado (la regresión de cuantiles con LightGBM en modo poblacional, con
XGBoost como alternativa equivalente) sobre el conjunto de prueba: para cada par
(ave, día) un punto $(\mathrm{lat}_t + p_{50,\mathrm{lat}}, \mathrm{lon}_t +
p_{50,\mathrm{lon}})$, un intervalo $[p_{10}, p_{90}]$ por componente con cobertura
empírica próxima al 80\,\% nominal, y la etiqueta de régimen del HMM causal del
capítulo \ref{ch:o3}. Los mapas del capítulo siguiente cartografían sobre esos tres
elementos el error de la predicción y lo estratifican por régimen y por mes, lo que
cierra la cadena Markov $\to$ HMM $\to$ predicción supervisada con la
visualización del comportamiento del modelo allí donde el error existe y allí donde
no.
